\section{Discussion and Conclusion}\label{sec:discussion}

Regional allocation and coverage are distinct properties of this post-hoc
Gaussian. Likelihood and scales determine covariance allocation; the predictive
centre also changes the error target and coverage. The tested training-only
choices do not consistently calibrate both regions.

\emph{Predictive centre and linearisation.}
The baseline mismatch survives removal of zero-mean prior shrinkage. At weaker
precision, affine prediction can improve without regional calibration, and a
modest relative parameter shift can produce a large nonlinear remainder.

Held-out refits demonstrate feature expressivity, separately from coefficient
selection. Spectral weights explain scale dependence, but the conditional limits
do not bound intermediate contrast, and coarse envelopes cannot certify the
baseline deficit. Evidence boundaries and perturbation sensitivity leave
weak-precision scale selection unresolved.

\emph{Residual-informed likelihoods.}
Residual-informed or block-wise variances $D$ replace $F$ by $J^\top D^{-1}J$,
after factoring out global noise variance, changing spectral directions and
within-region allocation. The common-scale oracle does not test this model. The
rejected Kuramoto--Sivashinsky example links small aggregate residuals to large
solution errors, without ruling out useful local residual information. Such
likelihood design remains untested.

\emph{Limitations.}
The study covers two $1{+}1$-dimensional nonlinear equations and two architectures with
a shared sub-network/fusion-head design. Only Burgers--MLP supplies eight fully
admitted seeds; the remaining families are small selected samples. Admission thresholds,
raw residual weights and reference-based regions are experimental conventions.
Root-coverage envelopes concern the five observed trained-centre crossings and
the stated bulk calibration, rather than optimal joint calibration.

The $\theta$-centred prior is conditional on fitted parameters. Nonlinear forward
checks retain covariance at the original checkpoint; they do not establish a nonlinear
posterior approximation. Spectral calculations do not certify weak-mode forward accuracy
or exact rank. The synthetic
stress test exceeds its whole-grid evidence tolerance, despite agreement at selected
settings (Appendix~\ref{app:r2}). Held-out expressivity is last-layer and within-grid;
it does not establish off-grid or full-network held-out approximation. Transfer to
higher dimensions, other architectures and training objectives is untested.

\emph{Implications.}
Assess accuracy, regional contrast and coverage for each predictive centre;
ensemble sensitivity likewise warrants assessment of the particular ensemble.
The trained-centre Gaussian is a covariance diagnostic.
An archived random half-grid split-conformal illustration \citep{lei2018conformal} on the reference run,
using absolute-error scores at target $0.95$, gives marginal coverage $0.950$ but
bulk, transition and sharp coverage $0.98$, $0.99$ and $0.31$. This descriptive
finite-grid example shows how marginal coverage can conceal regional imbalance.

\emph{Code and data availability.}
Code, trained parameters, saved results and the Lean development are available at
\url{https://github.com/soumya1729/pinn-uq-scale-sensitivity}.
The tagged reproduction release supplies the saved-data archive and its checksum.

\emph{Conclusion.}
Under the specified likelihood, sharp intervals under-cover at all five trained-centre
roots despite bulk calibration and regional RMS contrast agreement. Allocation
and coverage therefore require separate regional checks for each run and
predictive centre.
